\section{Introduction}

\headingindent
Modern image classifiers assume that test inputs resemble the data used for
training, these inputs are called In-Distribution (ID). That assumption fails in open environments. A medical classifier may
receive an image from another acquisition protocol, and a driving system may
encounter an object or visual condition absent from its training set. These abnormal inputs are called Out-Of-Distribution (OOD). A network
can then return a confident prediction even though none of its classes provides
a valid description of the input. Out-of-distribution detection adds a scalar
score and a rejection rule to identify such cases before the prediction is
used.

This internship started from a proposal to use model distillation as an
OOD detector, since this area is yet poorly explored in OOD detection related scientific papers. A high-capacity teacher is fixed after training on in-distribution
data. A smaller student learns only from those data and tries to reproduce a
teacher output or representation. If the student captures regularities that
hold on the training distribution, teacher-student disagreement may increase
on OOD inputs. 
% The original internship brief proposed this disagreement as an
% uncertainty score and set three tasks: implement the student training pipeline,
% build a reproducible evaluation protocol, and compare the resulting scores
% with established baselines.

The project also explored two extensions beyond matching the teacher's class
outputs. In the subspace-student ensemble, each student is trained on a fixed
or randomly sampled subset of an internal teacher representation. Disagreement
across the ensemble measures how predictions change when each student receives
only part of the representation. Feature denoising follows a different
pipeline and diverges from classical knowledge distillation: a student
receives a corrupted teacher feature and learns to reconstruct its clean
version. This direction was partly inspired by LeCun's proposal for predictive
representation learning and by joint-embedding predictive architectures
(JEPA), which predict representations from incomplete context rather than
reconstructing raw pixels \cite{lecun2022path,assran2023ijepa}.

\subsection{Progression of OOD detection methods}

\headingindent
% The papers named in the internship brief trace a change in where OOD evidence
% is sought.
The exploration of methods to detect OOD inputs has been in progress since 2017. Early work used the final softmax output. Later methods measured
distance in representation space, changed the training distribution, or
modified internal activations at inference time.

\begin{table}[ht]
\centering
\rowcolors{2}{rowblue!35}{white}
\begin{tabularx}{\textwidth}{@{}L{1.3cm}L{3.0cm}Y@{}}
\toprule
Year & Method & Main idea \\
\midrule
2017 & Maximum softmax probability (MSP) \cite{hendrycks2017baseline} &
Use the largest predicted class probability as an ID-confidence score. The
method needs no retraining, but a classifier can remain confident on OOD data. \\
2018 & Mahalanobis distance \cite{lee2018mahalanobis} &
Fit class-conditional Gaussian statistics to features from several network
layers. Score a test point by its distance to the closest class mean. \\
2021 & ReAct \cite{sun2021react} &
Clip unusually large penultimate activations before computing a confidence
score. The clipping threshold comes from ID data and reduces OOD
overconfidence. \\
2022 & Deep $k$-NN \cite{sun2022knn} &
Normalize deep features and measure the distance to nearby ID training
features. This removes the Gaussian assumption of Mahalanobis scoring. \\
2022 & PixMix \cite{hendrycks2022pixmix} &
Train with random compositions of images and structurally complex mixing
images. The method changes training data to improve several safety measures,
including anomaly detection. \\
\bottomrule
\end{tabularx}
\caption{A short timeline of the methods that motivated the internship.}
\label{tab:ood-timeline}
\end{table}

These methods served as inspiration for some of the experimental axes explored
here. MSP is a teacher-only baseline. Mahalanobis and $k$-NN motivate representation-space
scores. ReAct motivates activation clipping. PixMix motivates pixel
perturbation and feature prediction under image corruption.
% The
% internship asked whether a student trained only on ID data could turn these
% ideas into a family of teacher-student OOD Scores.

\subsection{Research questions}

\headingindent
The work addressed the following questions.

\begin{enumerate}
    \item Does a small student disagree with its teacher more often on OOD
    inputs than on ID inputs?
    \item Does corruption during student training expose a stronger signal than clean distillation?
    \item Should students reproduce class outputs, reconstruct internal
    features, or form an ensemble trained on different feature subspaces?
    \item How stable are the conclusions across teacher architectures, ID
    datasets, OOD datasets, score definitions, and evaluation conventions?
\end{enumerate}

% The last question became central. The final experiments showed that benchmark
% version and metric polarity can change an apparent comparison even when the
% underlying scores are unchanged.
